\documentclass[11pt]{article}
\usepackage[utf8]{inputenc}
\usepackage{float}
\usepackage{amsmath}
\newcommand{\mathbbm}[1]{\text{\usefont{U}{bbm}{m}{n}#1}} % from mathbbm.sty
\usepackage{amssymb}
\usepackage{soul} %to hilghlight with \hl{} and to strikeout with \st{}
\usepackage{color}  %to hilghlight with color with \hl{}
\usepackage{mathtools}
\usepackage{graphicx}
\usepackage{indentfirst}
\usepackage{enumitem}
\usepackage{setspace} 
%\setstretch{1.5}
\onehalfspacing
%\doublespacing
\usepackage{float}
\usepackage{fullpage}
\usepackage{ragged2e}
\usepackage[official]{eurosym}
\usepackage{booktabs}
\usepackage[bottom]{footmisc}
\newcommand{\sym}[1]{{#1}} % for symbols in Table
\usepackage{comment}
\usepackage{parskip}
\usepackage{tikz}
\usepackage[T1]{fontenc}
\usepackage{enumitem}
\setlist[itemize]{noitemsep}
\usepackage{placeins}
\usepackage[left=1in,right=1in,top=1in,bottom=1in]{geometry}


\usepackage{natbib}
\newcommand{\customcite}[1]{\citeauthor{#1}, \citeyear{#1}}
\bibliographystyle{chicago}
\usepackage{titlesec}
\usepackage[font=scriptsize]{subcaption}
	\captionsetup[figure]{font={small},labelfont={sc}}
	\captionsetup[table]{font={small},labelfont={sc}}
\titleformat*{\section}{\Large\bfseries}
\setcitestyle{authoryear,open={(},close={)}}
\usepackage{xcolor}
\usepackage[colorlinks]{hyperref}
\hypersetup{allcolors=blue!80!white}
\usepackage[scaled]{helvet}
\renewcommand{\familydefault}{\sfdefault}
 
\begin{document}
 
\begin{center}
    \textbf{\large \scshape CRISTIAN ARCIDIACONO}
\end{center}

% TODO: consider adding your office address here, e.g.
% Department of Economics, University of Bern, Bern, Switzerland
\begin{center}
    {\href{mailto:cristian.arcidiacono@unibe.ch}{cristian.arcidiacono@unibe.ch} $\vert$ \href{https://sites.google.com/view/cristian-arcidiacono/home}{Personal Page}}
\end{center}
\vspace{-0.2cm}
\noindent\rule{\textwidth}{0.4pt}

\vspace{0.2cm}

\begin{center}
    {\large \scshape Cover Letter}
\end{center}

\vspace{0.3cm}

Dear Members of the Search Committee,

I am a Ph.D. candidate in the joint Economics and Macro-Finance program at the University of Bern, supervised by Professor Pierpaolo Benigno. While my CV lists what I have done, this letter is about what I would bring as [POSITION] at [INSTITUTION]: a research agenda on inflation and the scarcity of safe assets, a network that spans academic and policy institutions and a path into economics that began outside it.

My research explores what drives inflation and what makes bonds special, and one question connects the two: how does scarcity in individual markets interact with conditions shared across the economy? My job market paper, \textit{Labor Market Shortages and the Phillips Curve: Identification under Heterogeneous Exposure to Common Shocks}, investigates whether the cross-sectional designs used to estimate the Phillips curve identify its slope. I show that they do not when units differ in their exposure to common shocks: the estimate conflates the response to a unit's own tightness with the response to tightness shared across the economy. I develop a two-step extension of interactive fixed effects that separates the two under testable conditions and apply it to firms' reported labor shortages across euro area sectors and countries. Inflation responds about five times more strongly to synchronized shortages than to micro-local ones, and the usual design recovers only a third of the local response. The paper received the 2026 INFER Best PhD Paper Award and I will present it at the ASSA Annual Meeting in January 2027.

A second strand of my research focuses on to the scarcity of safe assets. \textit{Dangerous Liaisons? Debt Supply and Convenience Yield Spillovers in the Euro Area}, with Matthieu Bellon and Matthias Gnewuch, asks what happens to a sovereign bond's convenience yield when other countries issue more debt. Using high-frequency identification around debt management announcements, we find that German and French issuance lowers convenience yields across the euro area, one-for-one for similarly safe issuers and less for riskier ones, as a model in which safe bonds are close substitutes predicts. The paper is accepted at the \textit{Review of Finance} and received the 2024 ECMI Best Paper Award.

[FIT for the INSTITUTION -- for places with strong interest; brief, one paragraph]

Beyond my research, I bring a network that spans academic and policy institutions. Over the past four years I have worked in six cities and four countries, the last of them London, where I spent the 2025--26 academic year visiting the LSE, hosted by Professor Ricardo Reis. Discussions there led me to the two-step design that now anchors the job market paper. Instead, earlier research positions at the Bank for International Settlements and the European Stability Mechanism taught me how a question changes when an institution must act on its answer, and the economists I met have become a network I rely on for feedback, discussion and co-authorship.

I also bring the perspective of an economist shaped by more than seminar rooms. I have been at the Military Academy of Modena, one of 95 admitted from over 10,000 applicants and after it, while completing my bachelor's degree, I ran the customer care department of a Milanese start-up for three years, and that job funded my master's in Geneva. In those same years, an internship at BlackRock taught me how to present an idea to people whose decisions depend on it, a style you will recognize in my job market talk. Through these experiences, I developed a sense of ownership over what I deliver and I learned to work under pressure and turn ideas into results; I have carried all of that into research. % They also gave me my questions. I financed my studies at a time when my family's finances were under strain, and I know what rising prices and debt mean for a household, not only from the data. 

Finally, I have been teaching online for nearly three years, a format that universities are increasingly developing. At UniDistance Suisse, a distance-learning university, I teach exercise sessions in French in ``International Finance'' and ``Money and Finance'', and supervise bachelor's theses. Without a physical classroom, I have learned to plan each session carefully, prepare materials students can revisit, and check their understanding regularly. Many of my students work alongside their studies or entered university through non-traditional paths. As a first-generation student who worked through my own degrees I relate to them and, helping them progress, is among the most rewarding parts of my work. %At the Geneva Graduate Institute I also co-supervised five applied research projects with twenty master's students.

Thank you for considering my application. I would welcome the opportunity to discuss my research and what I could contribute to [INSTITUTION]. \\

Sincerely,

Cristian Arcidiacono
















\begin{comment}
Dear Members of the Search Committee,

I am a Ph.D. candidate in Economics at the University of Bern, supervised by Professor Pierpaolo Benigno, and I expect to complete my degree in May 2027. I am applying for the [POSITION] at [INSTITUTION]. My research is in macroeconomics and finance, and a common question connects my research: how does scarcity in individual markets interact with conditions shared across the economy? I study how labor shortages affect inflation when they are local or widespread, and how debt issuance and collateral eligibility change the premia investors pay for safe assets.

My job market paper, \textit{Labor Market Shortages and the Phillips Curve: Identification under Heterogeneous Exposure to Common Shocks}, asks how labor scarcity affects inflation. Time fixed effects do not fully absorb common shocks when their effects differ across units. The estimated slope then blends the response to a sector's own tightness with the response to tightness shared across the economy. I develop a two-step extension of interactive fixed effects that separates the two under testable conditions. Using firms' reported labor shortages across euro area countries and manufacturing sectors, I find that a one standard deviation rise in a sector's own tightness raises its inflation by about 0.25 percentage points, three times the time fixed effects estimate, while a one standard deviation rise in synchronized tightness raises it by about 1.3 percentage points. The Phillips curve also steepens at high levels of tightness, and the non-linearity is concentrated in the synchronized component. The paper received the 2026 INFER Best PhD Paper Award, and I will present it at the ASSA Annual Meeting in January 2027.

A second strand of my research examines how debt supply affects the scarcity premium on safe bonds. In \textit{Dangerous Liaisons? Debt Supply and Convenience Yield Spillovers in the Euro Area}, with Matthieu Bellon and Matthias Gnewuch, we use high-frequency identification around debt management announcements to show that increases in German and French debt supply lower convenience yields beyond national borders. Spillovers to low-risk issuers are one-for-one and smaller for riskier ones, consistent with a model in which safe sovereign bonds are close substitutes. The paper is forthcoming in the \textit{Review of Finance} and received the 2024 ECMI Best Paper Award. Related work with Boris Hofmann, Max Riedel and Mathias Skrutkowski isolates the premium investors pay for Eurosystem collateral eligibility in corporate bond markets. In related work on inflation, a project with Boris Hofmann documents that, across 130 years and 17 advanced economies, faster productivity growth is associated with lower inflation and unchanged long-term nominal rates, a historical benchmark for the current debate on AI.

[FIT for the INSTITUTION]

% ETH Zurich example (one intellectual connection, one direction for collaboration):
% My job market paper builds on the European Commission business surveys, which connects directly to KOF's expertise in business tendency surveys, and I would welcome the opportunity to extend this work with KOF's survey data. A second direction for collaboration is my ongoing work with Boris Hofmann on productivity and inflation, which I plan to take to firm-level data with colleagues at the Chair of Applied Macroeconomics and at KOF.

My training has combined academic and policy environments. I spent the 2025--26 academic year at the LSE as a visiting scholar hosted by Professor Ricardo Reis, a year of seminars and presentations that showed me where the frontier of the field is and how I can contribute to it. Earlier positions at the Bank for International Settlements and the European Stability Mechanism taught me how a question changes when an institution must act on its answer, and each led to one of the coauthored projects above. Before my Ph.D., roles at BlackRock and AlphaMundi showed me how investors read economic developments, while leading a platoon of thirty soldiers at the Military Academy of Modena and running the customer care department of a Milanese start-up for three years taught me to meet deadlines and take responsibility for what a team delivers.

As a first-generation student, I financed my own studies, and that experience shapes my teaching. At UniDistance Suisse, I teach exercise sessions in French in ``International Finance'' and ``Money and Finance'' and supervise bachelor's theses. Many of my students work alongside their studies or entered university through non-traditional paths, and helping them progress is among the most rewarding parts of my work. Teaching online has taught me to structure sessions carefully, prepare materials that support independent study and regularly check students' understanding. At the Geneva Graduate Institute, I also co-supervised five applied research projects with twenty master's students.

Thank you for considering my application. I would welcome the opportunity to discuss my research and how I could contribute to [INSTITUTION]. \\

Sincerely,

Cristian Arcidiacono
\end{comment}

\end{document}
